%%%%%%%%%%%%%%%%%%%%%%%%
\chapter{Introdução}\label{intro}
%%%%%%%%%%%%%%%%%%%%%%%%

Em diversos problemas nos quais se utilizam técnicas de aprendizado de máquina, é importante que os modelos gerados sejam interpretáveis, uma vez que, além das predições realizadas, deseja-se compreender o conhecimento extraído dos dados. Essa característica é particularmente relevante em aplicações nas quais os resultados precisam ser analisados ou validados por especialistas. Dentre os modelos que favorecem a interpretação encontram-se aqueles baseados em regras, nos quais o conhecimento é representado por meio de condições que podem ser examinadas para compreender como as predições são realizadas \cite{Rudin2019,margot2021interpretability,molnar2022}.

Dentre os modelos baseados em regras encontram-se os classificadores associativos (CAs), que combinam as tarefas de mineração de regras de associação e classificação supervisionada \cite{liu1998}. Na classificação associativa (CA), são utilizadas regras de associação de classe (RACs), representadas na forma $A \Rightarrow c$, em que $A$ corresponde a um conjunto de condições presentes no antecedente e $c$ a uma classe. Essas regras são utilizadas na construção de modelos capazes de predizer a classe de novas instâncias. Por apresentarem explicitamente as condições utilizadas nas predições, as regras podem ser inspecionadas e analisadas por especialistas.

Ao longo dos anos, diferentes algoritmos de CA foram propostos, os quais adotam estratégias distintas para a geração, a ordenação (ranqueamento) e a poda das regras \cite{10.1016/j.eswa.2025.127454}. Embora a representação do conhecimento por meio de regras favoreça a interpretação, dependendo do número de regras e de condições presentes nas mesmas, essa tarefa pode não ser simples. Um modelo composto por centenas ou milhares de regras, por exemplo, dificilmente poderá ser analisado em sua totalidade por um especialista. Além disso, antecedentes extensos e regras redundantes também podem dificultar sua compreensão. Assim, além do desempenho preditivo, é importante analisar a quantidade e a estrutura das regras que compõem os modelos.

Esse problema torna-se especialmente importante quando se trabalha com conjuntos de dados de alta dimensionalidade. Na CA, cada atributo pode originar diferentes itens, correspondentes aos seus possíveis valores. Dessa forma, o aumento do número de atributos, especialmente daqueles que apresentam muitos valores distintos, amplia o conjunto de itens disponíveis para a extração das regras. Consequentemente, o espaço de busca pode aumentar, tornando o processo de indução mais custoso. Além disso, atributos irrelevantes ou redundantes podem contribuir para a geração de regras pouco informativas, aumentando a quantidade e a complexidade das regras produzidas.

Uma alternativa para tratar esse problema consiste em utilizar métodos de seleção de atributos (\textit{feature selection} -- FS). Esses métodos buscam identificar um subconjunto dos atributos originais, removendo aqueles considerados irrelevantes ou redundantes para a tarefa de aprendizado \cite{Dhal2022,Theng2024}. Diferentemente dos métodos de transformação, a FS preserva a representação original dos atributos selecionados, mantendo seu significado. Essa característica é interessante para modelos interpretáveis, como os classificadores associativos. No contexto da CA, a seleção de atributos pode influenciar não apenas o desempenho preditivo, mas também a quantidade e a estrutura das regras extraídas. Portanto, sua avaliação deve considerar tanto a redução de atributos quanto os efeitos dessa redução sobre a complexidade dos classificadores induzidos.

Diferentes estratégias para incorporar FS à CA encontram-se disponíveis na literatura (vide Capítulo~\ref{trabRelac}). Parte dos trabalhos realiza a seleção como uma etapa de pré-processamento, de modo que apenas os atributos selecionados sejam utilizados posteriormente na indução dos classificadores. Outros integram a seleção ao próprio processo de classificação, utilizando medidas de relevância durante a geração, a ponderação ou a avaliação das regras. Embora existam trabalhos que exploram essa integração, a combinação de diferentes métodos de seleção de atributos no contexto da CA ainda foi relativamente pouco investigada.

Diversos métodos de FS encontram-se disponíveis na literatura, os quais diferem entre si em relação aos critérios utilizados para avaliar a relevância dos atributos \cite{KHAIRE20221060,Theng2024}. Alguns utilizam medidas de associação estatística, enquanto outros consideram informação, correlação, dependência, redundância ou interação entre atributos. Dependendo do critério adotado e das características do conjunto de dados, diferentes métodos podem selecionar subconjuntos distintos. Além disso, não existe um método de FS que seja superior aos demais em todos os problemas. Assim, a escolha de um único seletor pode tornar o processo dependente de suas características e limitações.

Uma alternativa para diminuir essa dependência consiste em utilizar diferentes métodos de FS de maneira conjunta, por meio de um \textit{ensemble feature selection} (EFS). Nesse tipo de abordagem, os resultados de múltiplos processos de seleção são combinados para obter uma decisão sobre os atributos que deverão ser mantidos. A diversidade entre os componentes do \textit{ensemble} pode ser obtida, por exemplo, por meio de perturbações dos dados, diferentes configurações de um mesmo método ou pela combinação de seletores heterogêneos, que utilizam critérios distintos de avaliação dos atributos \cite{bolon2018scalability,mera2021framework}.

A utilização de métodos heterogêneos permite considerar diferentes perspectivas sobre a relevância, a redundância e a dependência entre os atributos. Entretanto, o uso de vários seletores não determina, por si só, como seus resultados devem ser combinados. A construção de um EFS envolve diferentes decisões metodológicas, dentre elas: (i) quais métodos utilizar; (ii) como representar e comparar os resultados produzidos; (iii) como analisar a diversidade entre os seletores; (iv) como combinar suas saídas; e (v) como definir o subconjunto final de atributos. Essas decisões podem influenciar tanto os atributos selecionados quanto os resultados obtidos pelos classificadores.

Quando os métodos de FS produzem ordenações dos atributos, seus resultados podem ser representados por meio de \textit{rankings}. Nesse caso, diferentes estratégias de agregação podem ser utilizadas para obter uma ordenação consensual. A literatura apresenta métodos eleitorais, técnicas de tomada de decisão multicritério, operadores de agregação, mecanismos de ponderação e abordagens de otimização \cite{drotar2019ensemble,hashemi2021pareto,abasabadi2021automatic,abellana2022ensemble,kaur2023mcdm}. Dentre essas alternativas encontram-se o Borda, que combina as posições atribuídas aos atributos pelos diferentes métodos, e o TOPSIS, que permite tratar a agregação como um problema de decisão multicritério, no qual os atributos são considerados alternativas avaliadas segundo diferentes critérios.

Além da agregação, é necessário definir quais atributos do \textit{ranking} consensual serão efetivamente selecionados. Embora o \textit{ranking} estabeleça uma ordem de preferência, ele não determina, necessariamente, quantos atributos devem ser mantidos. Para tanto, podem ser utilizadas estratégias baseadas em quantidades fixas, percentuais de retenção ou mecanismos adaptativos \cite{seijopardo2019automatic}. Dependendo da estratégia adotada, diferentes subconjuntos podem ser obtidos, com possíveis efeitos sobre o desempenho preditivo e a complexidade dos classificadores. Assim, a definição do subconjunto final constitui outra decisão importante na construção de um EFS.

Embora diferentes estratégias de EFS tenham sido propostas na literatura, sua aplicação à CA ainda é relativamente pouco explorada (vide Capítulo~\ref{trabRelac}). Os trabalhos identificados incluem tanto abordagens que utilizam métodos individuais de FS quanto propostas que combinam diferentes seletores antes da indução dos classificadores associativos. Dentre essas propostas destaca-se o trabalho de \citeonline{sowan2023hybrid}, que investiga uma configuração específica de \textit{ensemble}. Contudo, ainda são necessárias investigações sobre como diferentes decisões envolvidas na construção de um EFS heterogêneo, como a composição do \textit{ensemble}, a diversidade entre os ranqueadores, a estratégia de agregação e a definição do subconjunto final, afetam o desempenho preditivo e a complexidade dos classificadores associativos.

No contexto da CA, essa análise é importante porque diferentes estratégias de seleção de atributos podem produzir classificadores com desempenho preditivo semelhante, mas com quantidades e estruturas distintas de regras. Desse modo, não basta verificar se a seleção mantém ou melhora o desempenho preditivo. É necessário analisar também a redução de atributos e seus efeitos sobre a complexidade dos modelos induzidos.

Diante do exposto, neste trabalho propõe-se o EFS4AC (\textit{Ensemble Feature Selection for Associative Classification}), um \textit{pipeline} de EFS baseado na combinação e agregação de \textit{rankings} produzidos por diferentes métodos de seleção de atributos. A proposta utiliza métodos heterogêneos de FS e contempla cinco etapas principais, a saber: (i) seleção dos métodos que irão compor o \textit{ensemble}; (ii) geração e padronização dos \textit{rankings} individuais; (iii) análise da diversidade entre os \textit{rankings}; (iv) agregação dos resultados para obtenção de um \textit{ranking} consensual; e (v) definição do subconjunto final de atributos. O EFS4AC é utilizado na etapa de pré-processamento dos dados, antes da indução dos CAs, sem que seja necessário modificar seus algoritmos.

A hipótese deste trabalho é que a combinação de diferentes métodos de seleção de atributos, por meio de um \textit{pipeline} de EFS baseado na agregação de \textit{rankings}, possa reduzir a dependência de um único critério de avaliação dos atributos e produzir subconjuntos capazes de manter ou melhorar o desempenho preditivo dos classificadores associativos, ao mesmo tempo em que reduzem a complexidade dos modelos gerados. Essa hipótese não pressupõe que o uso do \textit{ensemble} seja vantajoso em todos os problemas, mas busca investigar se a combinação de diferentes critérios de relevância permite obter um melhor equilíbrio entre desempenho preditivo e complexidade.

Assim, busca-se responder à seguinte questão de pesquisa: \textit{Como diferentes decisões de projeto de um pipeline de EFS orientado à classificação associativa influenciam o desempenho preditivo e a complexidade dos modelos induzidos? Em quais condições sua utilização proporciona um melhor equilíbrio entre esses aspectos, quando comparada ao uso individual dos métodos de seleção de atributos e à classificação associativa sem seleção?}

Para tanto, tem-se como objetivo deste trabalho propor e avaliar o EFS4AC, um \textit{pipeline} integrado e configurável de \textit{Ensemble Feature Selection} orientado à classificação associativa. Busca-se investigar diferentes estratégias para a composição do \textit{ensemble}, a análise da diversidade entre os seletores, a agregação dos \textit{rankings} e a definição do subconjunto final de atributos. A avaliação será realizada considerando o desempenho preditivo, a redução de atributos e a complexidade dos modelos gerados, comparando-se as diferentes configurações do EFS4AC com métodos individuais de FS e com classificadores associativos sem seleção de atributos.

Como \textbf{contribuição metodológica} deste trabalho tem-se a proposta do EFS4AC, um \textit{pipeline} integrado e configurável de \textit{ensemble} de seleção de atributos voltado à classificação associativa. Além da obtenção de um \textit{ranking} consensual, pretende-se explorar estratégias para a composição do \textit{ensemble}, considerando a diversidade entre os seletores, e para a definição do subconjunto final de atributos. A proposta considera as particularidades da CA, especialmente os efeitos da seleção sobre a indução e a complexidade dos modelos baseados em regras. Já como \textbf{contribuição experimental}, espera-se, a partir da comparação entre diferentes configurações do EFS4AC, métodos individuais de FS e classificadores associativos sem seleção, identificar as vantagens, limitações e condições de aplicabilidade da estratégia proposta. Para tanto, serão analisados o desempenho preditivo, a redução de atributos e a complexidade dos modelos, considerando a quantidade e a estrutura das regras geradas.

O restante desta dissertação encontra-se organizado da seguinte forma. O Capítulo~\ref{cap_fundamentos} apresenta os conceitos necessários ao entendimento deste trabalho. No Capítulo~\ref{trabRelac} são apresentados e discutidos os trabalhos relacionados. O EFS4AC é descrito no Capítulo~\ref{cap:Metodologia}, juntamente com as etapas e decisões metodológicas que compõem o \textit{pipeline} proposto. Os capítulos seguintes apresentam a avaliação experimental e as considerações finais.
